
\begin{frame}{Prezentare generală}
\tableofcontents
\end{frame}

%-------------------------------------------------------
\section{Ce încercăm să rezolvăm?}
%-------------------------------------------------------

%-------------------------------------------------------
\begin{frame}{Verificarea limbajelor imperative}

\begin{itemize}
	\item Limbajele imperative sunt construite în jurul ideii de \alert{stare a programului} \\(date stocate în memorie).
	\medskip
	\item Programele imperative sunt secvențe de \alert{comenzi care modifică această stare}.
	\medskip \pause
	\item Pentru a \intens{demonstra proprietăți ale programelor imperative}, avem nevoie de:
	\begin{enumerate}
		\item O modalitate de a scrie afirmații despre starea programului.
		\item Reguli pentru manipularea și demonstrarea acestor afirmații.
	\end{enumerate}
	\medskip \pause
	\item Acestea vor fi oferite de \alert{Logica Hoare}.
\end{itemize}

\end{frame}
%-------------------------------------------------------

%-------------------------------------------------------
\section{Logica Hoare}
%-------------------------------------------------------

%-------------------------------------------------------
\begin{frame}{C.A.R. (Tony) Hoare}

\alert{Logica Hoare} a fost introdusă de \intens{Tony Hoare}.

Tot el a inventat algoritmul \alert{Quicksort} în anul 1960 (avea 26 de ani).

\begin{center}
\includegraphics[scale=.5]{../images/Hoare.jpeg}
\end{center}

\end{frame}
%-------------------------------------------------------


\begin{comment}

%-------------------------------------------------------
\begin{frame}{Propositional Logic (PL)}

\begin{center}
\begin{tabular}{lclcl}
\alert{Syntax} && \alert{Semantics} && \alert{Calculus} \\ \hline

Atomic formulas && Truth values && Natural Deduction \\
(variables) && $\{true, false\}$  && \\

$p,q,r,\mydots$ &&  &&  \\ [1em] \hline

Connectives && Truth tables && Rules of inference \\
$\vee,\wedge,\to,\neg$ && && \\[1em] \hline

Formulas && Satisfiable &&  Proof\\
$\varphi,\psi,\mydots$ &&  &&  \\ [1em]

 && Validity &&  \\
\end{tabular}

\bigskip
\intens{Soundness and Completeness:} Formula $\varphi$ is provable iff $\varphi$ is valid
\end{center}
\end{frame}
%-------------------------------------------------------


%-------------------------------------------------------
\begin{frame}{Hoare Logic}

\begin{center}
\begin{tabular}{ccc}
\alert{Syntax} & \alert{Semantics} & \alert{Calculus} \\ \hline

$\neg \wedge \vee \to \forall \exists$ & FOL &  N/A\\ \hline

$= + - \geq \leq \mydots$ & Arithmetics &  N/A\\ \hline

\intens{\texttt{:= ; while}}  & State maps variables & N/A \\
\intens{\texttt{if then else}} & to values (no pointers) &  \\ \hline

\alert{$\{P\}$\texttt{ S }$\{Q\}$} &  If initial state satisfies $P$ & \sel{todo} \\
& and \texttt{S} terminates then & \\
& final state satisfies $Q$&

\end{tabular}
\end{center}
\end{frame}
%-------------------------------------------------------
\end{comment}

%-------------------------------------------------------
\begin{frame}{Triplete Hoare}

Un \alert{triplet Hoare}  \intens{$\{P\}$\texttt{ S }$\{Q\}$} are trei componente:
\begin{description}
	\item[P] o \alert{pre-condiție}
	\item[S] un \alert{fragment de cod}
	\item[Q] o \alert{post-condiție}
\end{description}

\alert{Pre-condiția} este o afirmație care se referă la \intens{starea de dinaintea} execuției codului.

\alert{Post-condiția} este o afirmație care se referă la \intens{starea de după} execuția codului.

\end{frame}
%-------------------------------------------------------

%-------------------------------------------------------
\begin{frame}{Sintaxă}

\alert{Pre-condiția} și \alert{post-condiția} se construiesc din \intens{variabilele programului}, \intens{numere}, \intens{relații aritmetice}, și vor folosi \intens{logică propozitională} pentru a combina afirmații simple.

\begin{example}
\begin{itemize}
	\item $x = 3$, $\neg (x = y)$, $\neg (0 \leq x)$
	\item $x = 4 \wedge y = 2$
	\item $(\neg y \leq x) \to (x = 2 * y)$
	\item $\texttt{true}$, $\texttt{false}$
\end{itemize}
\end{example}

Dacă avem nevoie de o logică mai puternică, putem folosi \alert{logica de ordinul I}.
\begin{example}
Pentru a exprima că valoarea variabilei de program $x$ este un pătrat perfect,
putem scrie $\exists n. x = n * n$.
\end{example}

\end{frame}
%-------------------------------------------------------

%-------------------------------------------------------
\begin{frame}{Semantica intuitivă}

Reamintim că o \alert{stare} este dată de valorile (întregi) ale variabilelor de program (locațiilor de memorie).

\medskip \pause
\alert{Triplet Hoare:} $\{P\}$\texttt{ S }$\{Q\}$
\begin{itemize}
	\item \intens{dacă} $P$ e adevărată în starea de dinaintea execuției lui $S$
	\item \intens{și}  execuția lui \texttt{S} se termină
	\item \intens{atunci} $Q$ e adevărată în starea de după execuția lui $S$
\end{itemize}


\end{frame}
%-------------------------------------------------------

%-------------------------------------------------------
\begin{frame}{Corectitudine parțială}

Logica Hoare exprimă \alert{corectitudine parțială}!
\begin{itemize}
	\item Un program este \alert{parțial corect} dacă dă răspunsul așteptat atunci când se termină.
	\item Nu greșește niciodată, dar e posibil să nu dea nici un răspuns.
\end{itemize}

\medskip \pause
\begin{example}
\vspace{-.2cm}
\begin{center}
$\{x=1\}$ \texttt{while x=1 do y:=2} $\{x=3\}$
\end{center}
este o afirmație \intens{adevărată} conform logicii Hoare.
\begin{itemize}
	\item dacă starea de dinaintea execuției satisface $x=1$ \alert{și} bucla while se termină, atunci starea de după satisface $x=3$. \sel{Dar bucla while nu se termină!}
\end{itemize}

\end{example}
\end{frame}
%-------------------------------------------------------

%-------------------------------------------------------
\begin{frame}{Corectitudinea parțială e suficientă}

\alert{De ce nu insistăm asupra terminării?}
\begin{itemize}
	\item Poate că nu vrem terminare.
	\item Simplifică logica (există variante și pentru corectitudine totală).
	\item Există metode specializate pentru demonstrarea terminării, dacă vrem.
\end{itemize}

% We will come back to termination with the \alert{Weakest Precondition Calculus}.
\end{frame}
%-------------------------------------------------------


%-------------------------------------------------------
\begin{frame}{Cum?}

\vspace{-.6cm}
\begin{example}
\smallskip
Tripletele Hoare ne permit să facem afirmații precum:
\vspace{-.2cm}
\begin{center}
$\{x>0\}$ \texttt{y:=0-x} $\{y < 0 \wedge x \neq y\}$
\end{center}
\pause
Dacă $(x > 0)$ este adevărată \intens{înainte} de execuția instrucțiunii \texttt{y:=0-x}
atunci $(y < 0 \wedge x \neq y)$ este adevărată \intens{după}.

\pause
Această afirmație este în mod intuitiv adevărată.
\intens{Cum o demonstrăm?}
\end{example}

\medskip \pause
%We can use pre- and postconditions to specify the effect of a code fragment on the state, but \intens{how do we prove or disprove a Hoare Triple specification}?

Avem nevoie de niște \intens{reguli} pentru a deduce (formal) astfel de triplete.
Vom avea câte o regulă pentru fiecare din cele cinci feluri de instrucțiuni (plus încă două alte reguli).
\end{frame}
%-------------------------------------------------------


%-------------------------------------------------------
\begin{frame}{Axioma pentru \texttt{skip} (Regula 0/6)}

\texttt{skip} nu face nimic, deci e de așteptat ca tripletele Hoare pentru \texttt{skip} reflecte acest lucru.

\medskip
\alert{Axioma pentru \texttt{skip}:}
\begin{center}
\fbox{\intens{$\{Q\}$ \texttt{skip} $\{Q\}$}}
\end{center}

\end{frame}
%-------------------------------------------------------


%-------------------------------------------------------
\begin{frame}{Axioma atribuirii (Regula 1/6)}

Atribuirile \sel{schimbă starea}, deci e de așteptat ca tripletele Hoare pentru atribuire să reflecte această schimbare.

\medskip \pause
\alert{Axioma atribuirii:}
\begin{center}
\fbox{\intens{$\{Q(e)\}$ \texttt{x:=e} $\{Q(x)\}$}}
\end{center}
{\footnotesize
($Q(x)$ este o proprietate asupra unei variabile $x$ și $Q(e)$ indică aceeași proprietate în care toate aparițiile lui $x$ au fost înlocuite de expresia $e$)
}

\medskip \pause
Dacă vrem ca $x$ să aibă o anumită proprietate $Q$ \sel{după} atribuire, atunci acea proprietate trebuie să țină pentru expresia $e$ atribuită lui $x$ \intens{înainte} de execuția atribuirii.
\end{frame}
%-------------------------------------------------------


%-------------------------------------------------------
\begin{frame}{Axioma atribuirii}

\begin{example}
Regula \false{inversă} este falsă: \false{$\{Q(x)\}$ \texttt{x:=e} $\{Q(e)\}$}

Dacă am aplica această "axiomă" greșită pre-condiției $x=0$ și fragmentului de cod \texttt{x:=1}, am obține
\begin{center}
\Hoare{x=0}{x:=1}{1=0}
\end{center}
care se citește "dacă $x=0$ initial și \texttt{x:=1} se temină, atunci $1=0$ după execuție".
\end{example}
\end{frame}
%-------------------------------------------------------

%-------------------------------------------------------
\begin{frame}{Lucrăm invers, dinspre post-condiție spre pre-condiție}

Ar putea părea natural să începem cu \false{pre-condiția} și să facem deducții \false{înspre post-condiție}, dar acesta nu e cel mai bun mod de a raționa folosind logica Hoare.

	\medskip
Dimpotrivă, \true{pornim cu post-condiția și mergem "înapoi"}.

\pause
	\begin{example}
		Pentru a aplica axioma pentru atribuire,
		\begin{center}
		\intens{$\{Q(e)\}$ \texttt{x:=e} $\{Q(x)\}$}
		\end{center}
		luăm post-condiția, o copiem în pre-condiție și apoi înlocuim toate aparițiile lui $x$ cu $e$.

		Observație: postcondiția poate avea una, mai multe, sau nici o apariție a lui $x$.
		\intens{Toate vor fi înlocuite cu $e$ în pre-condiție!}
	\end{example}

\end{frame}
%-------------------------------------------------------

%-------------------------------------------------------
\begin{frame}{Regula pentru atribuire}

\vspace{-.6cm}
\alert{Axioma atribuirii:} \intens{$\{Q(e)\}$ \texttt{x:=e} $\{Q(x)\}$}

\medskip
\begin{example}
Presupunem că avem de executat \intens{\texttt{x:=2}} și că \sel{post-condiția dorită} este \intens{$y=x$}.

O instanță a axiomei de atribuire:
\pause
\vspace{-.2cm}
\begin{center}
\Hoare{y=\intens{2}}{x:=2}{y=\intens{x}}
\end{center}
\end{example}
\end{frame}
%-------------------------------------------------------

%-------------------------------------------------------
%\begin{frame}{Regula pentru atribuire}
%
%You can always replace predicates by \alert{equivalent predicates}; just label your
%proof step with \intens{"precondition equivalence"} or \intens{"postcondition equivalence"}.
%
%\medskip \pause
%\begin{example}
%How should we prove
%\vspace{-.2cm}
%\begin{center}
%\Hoare{y>0}{x:=y+3}{x>3}?
%\end{center}
%\vspace{-.2cm} \pause
%Start with the postcondition $x>3$ and apply the assignment axiom:
%\vspace{-.2cm}
%\begin{center}
%\Hoare{\intens{y+3}>3}{x:=y+3}{\intens{x}>3}
%\end{center}
%\vspace{-.2cm} \pause
%Then use the fact that \intens{$y+3 > 3$} is equivalent with \intens{$y>0$} to get the result.
%\end{example}
%\end{frame}
%-------------------------------------------------------


%-------------------------------------------------------
\begin{frame}{Regula pentru atribuire}

\begin{example}
Să zicem că vrem să demonstrăm
\vspace{-.2cm}
\begin{center}
\Hoare{y=2}{x:=y}{x>0}?
\end{center}
\vspace{-.2cm}  \pause
Această aritmație este evident adevărată. Dar folosind axioma pentru atribuire obținem:
\vspace{-.2cm}
\begin{center}
\Hoare{\intens{y}>0}{x:=y}{\intens{x}>0}?
\end{center}
\vspace{-.2cm}
Și \ldots {\color{red} $y>0$ nu este echivalentă cu $y=2$!}
\end{example}
\end{frame}
%-------------------------------------------------------

%-------------------------------------------------------
\begin{frame}{Afirmații tari și slabe}

O afirmație $P$ este \alert{mai tare} decât $Q$ dacă \alert{$P$ o implică pe $Q$}.

Dacă $P$ este mai tare decât $Q$, atunci $P$ \intens{are mai multe șanse să fie falsă} decât $Q$.

\pause
\begin{example}[promisiune electorală]
\begin{itemize}
	\item {\em Voi menține șomajul sub $3\%$} este \alert{mai tare} decât
	\item {\em Voi menține șomajul sub $15\%$}
\end{itemize}
\end{example}

\pause
\alert{Cea mai tare} afirmație posibilă este \intens{$\bot$}.


\alert{Cea mai slabă} afirmație posibilă este \intens{$\top$}.

\end{frame}
%-------------------------------------------------------

%-------------------------------------------------------
\begin{frame}{Post-condiții tari}

\begin{example}
Triplul Hoare \Hoare{x=5}{x:=x+1}{\false{x=6}} este mai

informativ decât \Hoare{x=5}{x:=x+1}{\true{x>0}}.
\end{example}

\medskip \pause
Dacă \alert{post-condiția} $Q_1$ este \alert{mai tare} decât $Q_2$, atunci \\ \Hoare{P}{S}{Q_1} este o afirmație \alert{mai tare} decât \Hoare{P}{S}{Q_2}.

\medskip \pause
\begin{example}
Deoarece post-condiția \false{$x=6$} este \intens{mai tare} decât \true{$x>0$} (întrucât $x = 6 \to x > 0$),

atunci \Hoare{x=5}{x:=x+1}{\false{x=6}} este o afirmație \intens{mai tare}

decât \Hoare{x=5}{x:=x+1}{\true{x>0}}.
\end{example}
\end{frame}
%-------------------------------------------------------


%-------------------------------------------------------
\begin{frame}{Pre-condiții slabe}

\begin{example}
Tripletul Hoare \Hoare{\true{x>0}}{x:=x+1}{x>1} este mai

informativ decât \Hoare{\false{x=5}}{x:=x+1}{x>1}.
\end{example}

\medskip
Dacă o \alert{pre-condiție} $P_1$ este \alert{mai slabă} decât $P_2$, atunci \\ \Hoare{P_1}{S}{Q} este o afirmație \alert{mai tare} decât \Hoare{P_2}{S}{Q}.

\medskip
\begin{example}
Deoarece pre-condiția \true{$x>0$} este \intens{mai slabă} decât \false{$x=5$},

atunci \Hoare{\true{x>0}}{x:=x+1}{x>1} este o afirmație \intens{mai tare}

decât  \Hoare{\false{x=5}}{x:=x+1}{x>1}.
\end{example}
\end{frame}
%-------------------------------------------------------

%-------------------------------------------------------
\begin{frame}{Regula de întărire a pre-condiției (Regula 2/6)}

Este sigur (corect) să \alert{întărim} o \intens{pre-condiție}.

\alert{Regula de întărire a pre-condiției:}
\begin{center}
\fbox{\intens{$\infer[\mbox{ dacă }P_s \to P_w ]{\{P_s\}\ \mathtt{S}\ \{Q\}}{\{P_w\}\ \mathtt{S}\ \{Q\}}$}}
\end{center}

\begin{example}
\begin{center}
$\infer[\mbox{ deoarece }\false{y=2} \to \true{y>0}]{\false{\{y=2\}}\ \mathtt{x := y}\ \{x>0\}}{\true{\{y>0\}}\ \mathtt{x := y}\ \{x>0\}}$
\end{center}

\end{example}
\end{frame}
%-------------------------------------------------------

%-------------------------------------------------------
\begin{frame}{Regula de slăbire a post-condiției (Regula 3/6)}

Este sigur (corect) să {\em relaxăm} (\alert{slăbim}) o \intens{post-condiție}.

\alert{Regula de slăbire a post-condiției:}
\begin{center}
\fbox{\intens{$\infer[\mbox{ dacă } Q_s \to Q_w ]{\{P\}\ \mathtt{S}\ \{Q_w\}}{\{P\}\ \mathtt{S}\ \{Q_s\}}$}}
\end{center}

\begin{example}
\begin{center}
$\infer[\mbox{ deoarece }\false{x>3} \to \true{x>1}]{\{x>2\}\ \mathtt{x := x+1}\ \true{\{x>1\}}}{{\{x>2\}}\ \mathtt{x := x+1}\ \false{\{x>3\}}}$
\end{center}

\end{example}
\end{frame}
%-------------------------------------------------------

%-------------------------------------------------------
\begin{frame}{Regula pentru secvențiere (Regula 4/6)}

Programele imperative conțin instrucțiuni, care modifică starea \intens{secvențial}.

\alert{Regula de secvențiere:}
\begin{center}
\fbox{\intens{$\infer{\{P\} \mathtt{S_1 ; S_2} \{R\}}{\{P\} \mathtt{S_1} \{Q\} & \{Q\} \mathtt{S_2} \{R\}}$}}
\end{center}

\pause
\begin{example}
\begin{center}
$\infer{\HoareM{x>2}{\sel{x:=x+1\ ;\  x:=x+2}}{x>5}}{\HoareM{x>2}{\sel{x:=x+1}}{x>3} & \HoareM{x>3}{\sel{x:=x+2}}{x>5}}$
\end{center}

\end{example}
\end{frame}
%-------------------------------------------------------


%-------------------------------------------------------
\begin{frame}{Cum obținem condiția din mijloc?}

De obicei, când aplicăm o regulă de forma
\begin{center}
\fbox{\intens{$\infer{\{P\} \mathtt{S_1 ; S_2} \{R\}}{\{P\} \mathtt{S_1} \{Q\} & \{Q\} \mathtt{S_2} \{R\}}$}}
\end{center}
pre-condiția $P$ și post-condiția $R$ ne vor fi date; totuși, de unde obținem $Q$?

\pause
Pornind de la ținta noastră, $R$, și \intens{mergând înapoi}!

\begin{center}
$\infer{\HoareM{x>2}{\sel{x:=x+1 ; x:=x+2}}{x>5}}{\HoareM{x>2}{\sel{x:=x+1}}{Q} & \HoareM{Q}{\sel{x:=x+2}}{x>5}}$
\end{center}

\end{frame}
%-------------------------------------------------------

%-------------------------------------------------------
\begin{frame}{Cum scriem o demonstrație?}

\vspace{-.6cm}
\begin{example}
Să zicem că vrem să demonstrăm
\begin{center}
\Hoare{x=3}{\sel{x:=x+1; x:=x+2}}{x>5}.
\end{center}

\pause
\begin{enumerate}[<+->]
	\item \Hoare{x+2 > 5}{\sel{x:=x+2}}{x>5} \hfill (Atribuire)
	% \item \Hoare{x > 3}{\sel{x:=x+2}}{x>5} \hfill (1, Precondition Equivalence)
	\item \Hoare{(x+1) + 2 > 5}{\sel{x:=x+1}}{x+2 > 5} \hfill (Atribuire)
	% \item \Hoare{x > 2}{\sel{x:=x+1}}{x>3} \hfill (3, Precondition Equivalence)
	\item \Hoare{(x+1) + 2 > 5}{\sel{x:=x+1 ; x:=x+2}}{x>5} \hfill (1,2, Sevențiere)
	\item \Hoare{x=3}{\sel{x:=x+1; x:=x+2}}{x>5}. \hfill{(3, Întărirea pre-condiției)}

	deoarece $x=3 \to (x+1) + 2 > 5$
\end{enumerate}
\end{example}

Observați numerotarea pașilor și justificările!
\end{frame}
%-------------------------------------------------------

%-------------------------------------------------------
\begin{frame}{Regula pentru instrucțiunea {\tt if} (Regula 5/6)}

\alert{Regula pentru {\tt if}:}
\begin{center}
\fbox{\intens{$\infer{\{P\}\ \mathtt{if\ b\ then\ S_1\ else\  S_2}\ \{Q\}}{\{P \wedge b\}\ \mathtt{S_1}\ \{Q\} & \{P \wedge \neg b\}\ \mathtt{S_2}\ \{Q\}}$}}
\end{center}

\pause
\begin{itemize}
	\item Când se execută o instrucțiune {\tt if}, se va executa ori $\mathtt{S_1}$ ori $\mathtt{S_2}$.
	\item De aceea, ca în urma execuției să fie validă afirmația $Q$, trebuie ca \alert{ambele} ramuri $\mathtt{S_1}$ și $\mathtt{S_2}$ să facă $Q$ adevărată.
	\item Asemănător, dacă pre-condiția pentru {\tt if} este $P$, ea trebuie să fie \alert{o} pre-condiție pentru \alert{ambele} ramuri $\mathtt{S_1}$  și $\mathtt{S_2}$.
	\item Alegerea dintre $\mathtt{S_1}$ și $\mathtt{S_2}$ depinde de validitatea expresiei $b$ în starea inițială; deci putem presupune $b$ ca \alert{o} pre-condiție pentru $\mathtt{S_1}$ și $\neg b$ ca \alert{o} pre-condiție pentru $\mathtt{S_2}$.
\end{itemize}
\end{frame}
%-------------------------------------------------------

%-------------------------------------------------------
\begin{frame}{Regula pentru If}

\alert{Regula pentru {\tt if}:}
\begin{center}
\fbox{\intens{$\infer{\{P\}\ \mathtt{if\ b\ then\ S_1\ else\  S_2}\ \{Q\}}{\{P \wedge b\}\ \mathtt{S_1}\ \{Q\} & \{P \wedge \neg b\}\ \mathtt{S_2}\ \{Q\}}$}}
\end{center}

\begin{example}

Să zicem că vrem să demonstrăm

\begin{center}
\Hoare{x>3}{\sel{if  x>2 then y:=1 else y:=-1}}{y>0}
\end{center}
\pause
Regula pentru {\tt if} ne spune că e suficient să demonstrăm:
\begin{enumerate}[I]
	\item \Hoare{(x>3) \wedge \intens{(x > 2)}}{\sel{y:=1}}{y>0}
	\item \Hoare{(x>3) \wedge \intens{\neg(x>2)}}{\sel{y:=-1}}{y>0}
\end{enumerate}
\end{example}
\end{frame}
%-------------------------------------------------------

%-------------------------------------------------------
\begin{frame}{Regula pentru If}

\begin{example}[cont.]
Pentru I. \Hoare{(x>3) \wedge (x > 2)}{\sel{y:=1}}{y>0} avem demonstrația:
\begin{enumerate}[<+->]
	\item \Hoare{1>0}{y:=1}{y>0} \hfill (Atribuire)
	\item \Hoare{(x>3) \wedge (x > 2)}{y:=1}{y >0}  \hfill (1, Întărirea pre-condiției)

	deoarece $(x>3) \wedge (x > 2) \to 1>0$
\end{enumerate}

\medskip
\onslide<3->{Pentru II. \Hoare{(x>3) \wedge \neg(x>2)}{\sel{y:=-1}}{y>0} avem demonstrația:}
\begin{enumerate}[<+->]
{\setcounter{enumi}{2}}
	\item \Hoare{-1>0}{y:=-1}{y>0} \hfill (Atribuire)
	\item \Hoare{(x>3) \wedge \neg(x>2)}{y:=-1}{y>0}  \hfill (3, Întărirea pre-condiției)

	   deoarece $(x>3) \wedge \neg(x>2) \to -1 > 0 $
\end{enumerate}

\medskip
\onslide<5->{Putem concluziona:}
\begin{enumerate}[<+->]
	{\setcounter{enumi}{4}}
	\item \Hoare{x>3}{\sel{if  x>2 then y:=1 else y:=-1}}{y>0} \hfill (2, 4, If)
\end{enumerate}
\end{example}
\end{frame}
%-------------------------------------------------------

%-------------------------------------------------------
\begin{frame}{Regula pentru If}

{\color{red}\underline{Exercițiu:}}

Cum ați scrie o regulă pentru o instrucțiune condițională fără \intens{else}?
\begin{center}
\texttt{if b then S}
\end{center}

\end{frame}
%-------------------------------------------------------

%-------------------------------------------------------
\begin{frame}{Regula pentru While (Regula 6/6)}

\vspace{-.6cm}
\alert{Regula pentru While:}
\begin{center}
\fbox{\intens{$\infer{\{P\}\ \mathtt{while\ b\ do\ S}\ \{P\wedge \neg b\}}{\{P \wedge b\}\ \mathtt{S}\ \{P\}}$}}
\end{center}

\pause
\begin{itemize}
	\item $P$ este numit \alert{invariantul buclei}
	\item $P$ este adevărată înainte de a intra în buclă, și este păstrată adevărată de execuția corpului buclei, $S$ (deși nu neapărat și în timpul execuției lui $S$).
	\item Dacă execuția buclei se temină condiția de control trebuie să fie falsă, deci $\neg b$ apare în post-condiție.
	\item Pentru a executa corpul buclei, \texttt{S}, $b$ trebuie să fie adevărată, deci $b$ apare în pre-condiția lui $S$.
\end{itemize}
\end{frame}
%-------------------------------------------------------

%-------------------------------------------------------
\begin{frame}{Aplicarea regulii pentru While}

\begin{center}
\begin{tabular}{ccc}
$\infer{\{P\}\ \mathtt{while\ b\ do\ S}\ \{P\wedge \neg b\}}{\{P \wedge b\}\ \mathtt{S}\ \{P\}}$ &&  $\{P\}\ \mathtt{while\ b\ do\ S}\ \{Q\}$
\end{tabular}
\end{center}

\begin{itemize}
	\item \false{Partea cea mai grea} este descoperirea \alert{invariantului}.

	\item Este nevoie de \alert{intuiție}. \false{Nu există algoritm} care să găsească invariantul.

	\item Post-condiția obținută după aplicarea regulii e de forma $P\wedge \neg b$.

	 E posibil ca post-condiția dorită, $Q$, să fie diferită!

	\item Dacă $(P \wedge \neg b) \to Q$, putem folosi regula de \intens{slăbire a post-condiției}.
\end{itemize}

\begin{center}
$
\infer[\mbox{ deoarece }P \wedge \neg b \to Q]{\{P\}\ \mathtt{while\ b\ do\ S}\ \{Q\}}
{
\infer{\{P\}\ \mathtt{while\ b\ do\ S}\ \{P\wedge \neg b\}}{\{P \wedge b\}\ \mathtt{S}\ \{P\}}
}
$
\end{center}
\end{frame}
%-------------------------------------------------------

%-------------------------------------------------------
\begin{frame}{Regula pentru While}

\begin{example}
Să zicem că vrem să găsim o pre-condiție $P$ astfel încât
\begin{center}
\Hoare{P}{while (n>0) do n:=n-1}{n=0}
\end{center}
\pause
Vrem un invariant $P$ astfel încât
\begin{itemize}
	\item $P$ este menținut adevărat de corpul buclei
	\item $P \wedge \neg(n>0) \to (n=0)$
\end{itemize}
\pause
\alert{$P \equiv n\geq 0$} pare o alegere rezonabilă pentru un astfel de invariant.

Premiza regulii pentru while este satisfăcută de axioma de atribuire.


\end{example}
\end{frame}
%-------------------------------------------------------

%-------------------------------------------------------
\begin{frame}{Regula pentru While}

\begin{example}[cont.]
	Demonstrăm că
\begin{center}
\Hoare{n \geq 0}{while (n>0) do n:=n-1}{n=0}
\end{center}

\begin{enumerate}[<+->]
	\item \Hoare{n-1 \geq 0}{n:=n-1}{n \geq 0} \hfill (Atribuire)
	\item \Hoare{n\geq 0 \wedge n > 0}{n:=n-1}{n \geq 0} \hfill (1, Întărirea pre-condiției)

		deoarece $n\geq 0 \wedge n > 0 \to n-1 \geq 0$
	\item \Hoare{n \geq 0}{while (n>0) do n:=n-1}{n \geq 0 \wedge \neg(n > 0)} \hfill (2, While)
	\item \Hoare{n \geq 0}{while (n>0) do n:=n-1}{n = 0} \hfill (3, Slăbirea post-condiției)

		deoarece $n \geq 0 \wedge \neg(n > 0) \to n = 0$
\end{enumerate}
%\pause
%There are other invariants we could have chosen, such as $n=0$ and $\bot$.
%
%The one we chose was "better" because it gave us a \intens{weaker} precondition.
\end{example}
\end{frame}
%-------------------------------------------------------

%-------------------------------------------------------
\begin{frame}{Sistemul de reguli pentru Logica Hoare}

\begin{description}
\item[Skip:] \fbox{\intens{$\{Q\}$ \texttt{skip} $\{Q\}$}}
\vitem[Atribuire:] \fbox{\intens{$\{Q(e)\}$ \texttt{x:=e} $\{Q(x)\}$}}

\vitem[Întărire pre:] \fbox{\intens{$\infer[\mbox{ dacă } P_s \to P_w]{\{P_s\}\ \mathtt{S}\ \{Q\}}{\{P_w\}\ \mathtt{S}\ \{Q\}}$}}

\vitem[Slăbire post:] \fbox{\intens{$\infer[\mbox{ dacă } Q_s \to Q_w]{\{P\}\ \mathtt{S}\ \{Q_w\}}{\{P\}\ \mathtt{S}\ \{Q_s\}}$}}

\vitem[Secvențiere:] \fbox{\intens{$\infer{\{P\} \mathtt{S_1 ; S_2} \{R\}}{\{P\} \mathtt{S_1} \{Q\} & \{Q\} \mathtt{S_2} \{R\}}$}}

\vitem[If:] \fbox{\intens{$\infer{\{P\}\ \mathtt{if\ b\ then\ S_1\ else\  S_2}\ \{Q\}}{\{P \wedge b\}\ \mathtt{S_1}\ \{Q\} & \{P \wedge \neg b\}\ \mathtt{S_2}\ \{Q\}}$}}

\vitem[While:] \fbox{\intens{$\infer{\{P\}\ \mathtt{while\ b\ do\ S}\ \{P\wedge \neg b\}}{\{P \wedge b\}\ \mathtt{S}\ \{P\}}$}}
\end{description}
\end{frame}
%-------------------------------------------------------

%-------------------------------------------------------
\begin{frame}{Un program simplu}
\begin{example}
\intens{Suma primelor $n$ numere naturale impare este $n^2$.}

Program:

\hspace{.5cm} \texttt{i := 0;}

\hspace{.5cm} \texttt{s := 0;}

\hspace{.5cm} \texttt{while (i $\neq$ n) do}

\hspace{.8cm} \texttt{i := i+1;}

\hspace{.8cm} \texttt{s := s+(2*i-1)}

\medskip
Țintă: \Hoare{\top}{Program}{s = n^2}
\end{example}
\end{frame}
%-------------------------------------------------------

%-------------------------------------------------------
\begin{frame}{Un program simplu}
\begin{example}[cont.]
Să vedem câteva exemple:
\begin{itemize}
	\item $1 = 1 = 1^2$
	\item $1 + 3 = 4 = 2^2$
	\item $1 + 3 + 5 = 9 = 3^2$
	\item $1 + 3 + 5 + 7 = 16 = 4^2$
\end{itemize}

Pare să meargă. Să vedem dacă putem și demonstra!

\medskip
Tintă: \Hoare{\top}{Program}{s = n^2}
\end{example}
\end{frame}
%-------------------------------------------------------

%-------------------------------------------------------
\begin{frame}{Un program simplu}
\begin{example}[cont.]
\intens{Mai întăi, avem nevoie de un invariant $P$.}

\medskip
\begin{columns}
\begin{column}{0.4\textwidth}
\fbox{\intens{$\infer{\{P\}\ \mathtt{while\ \true{b}\ do\ \false{S}}\ \{P\wedge \neg \true{b}\}}{\{P \wedge \true{b}\}\ \mathtt{S}\ \{P\}}$}}
\end{column}
\begin{column}{0.4\textwidth}
\texttt{while \true{(i $\neq$ n)} do}

\hspace{.3cm} \false{\texttt{i := i+1;}}

\hspace{.3cm} \false{\texttt{s := s+(2*i-1)}}

$\{s = n^2\}$
\end{column}
\end{columns}

\medskip
Din regula pentru while, vrem ca \intens{$P \wedge (i = n) \to (s = n^2)$}, ca să putem aplica regula de slăbire a post-condiției.

La fiecare trecere prin corpul buclei, $i$ e incrementat și $s$ trece la următorul pătrat perfect.

\pause
\intens{Invariantul \alert{$P \equiv (s = i^2)$} pare rezonabil.}
\end{example}
\end{frame}
%-------------------------------------------------------


%-------------------------------------------------------
\begin{frame}{Un program simplu}
% \begin{example}[cont.]
Verificăm că $P \equiv (s = i^2)$ este invariant: \intens{deducem \Hoare{P \wedge (i \neq n)}{S}{P}}

\pause
\begin{enumerate}[<+->]
	\vitem \Hoare{s+(2*i-1)=i^2}{s:=s+(2*i-1)}{s=i^2} \hfill (Atribuire)
	\vitem \Hoare{s+(2*(i+1)-1)=(i+1)^2}{i:=i+1}{s+(2*i-1)=i^2} \hfill (Atribuire)
	\vitem \Hoare{s+(2*(i+1)-1)=(i+1)^2}{i:=i+1;\ s:=s+(2*i-1)}{s=i^2} \\ \hfill(1, 2, Sevențiere)
	\vitem \Hoare{s = i^2 \wedge i \neq n}{i:=i+1;\ s:=s+(2*i-1)}{s=i^2} \\ \hfill(3 Întărirea pre-condiției)

	  deoarece $s = i^2 \wedge i \neq n \to s+(2*(i+1)-1)=(i+1)^2$
\end{enumerate}
% \end{example}
\end{frame}
%-------------------------------------------------------

% \addtocounter{framenumber}{-1}
%-------------------------------------------------------
% \begin{frame}{A simple program}
% \begin{example}[cont.]
% Check $P \equiv (s = i^2)$ is an invariant: \intens{prove \Hoare{P \wedge (i \neq n)}{S}{P}}
%
% \begin{center}
% $\infer
% {\HoareM{s = i^2 \wedge i \neq n}{i:=i+1;\ s:=s+(2*i-1)}{s=i^2}}
% {\HoareM{s=i^2 \wedge i \neq n}{i:=i+1}{Q} & \HoareM{Q}{s:=s+(2*i-1)}{s=i^2}}$
% \end{center}
%
% $Q$ is $\{s+(2*i-1)=i^2\}$
% \begin{enumerate}
	% \item \Hoare{s+(2*i-1)=i^2}{s:=s+(2*i-1)}{s=i^2} \hfill (Assignment)
	% \item
	% \item \Hoare{s=i^2 \wedge i \neq n}{i:=i+1}{Q}
	% \item \Hoare{s = i^2 \wedge i \neq n}{i:=i+1;\ s:=s+(2*i-1)}{s=i^2} \hfill (1,3, Sequencing)
% \end{enumerate}
% \vspace{.65cm}
% \end{example}
% \end{frame}
% -------------------------------------------------------
%
% \addtocounter{framenumber}{-1}
% -------------------------------------------------------
% \begin{frame}{A simple program}
% \begin{example}[cont.]
% Check $P \equiv (s = i^2)$ is an invariant: \intens{prove \Hoare{P \wedge (i \neq n)}{S}{P}}
%
% \begin{center}
% $\infer
% {\HoareM{s = i^2 \wedge i \neq n}{i:=i+1;\ s:=s+(2*i-1)}{s=i^2}}
% {\HoareM{s=i^2 \wedge i \neq n}{i:=i+1}{Q} & \HoareM{Q}{s:=s+(2*i-1)}{s=i^2}}$
% \end{center}
%
% $Q$ is $\{s+(2*i-1)=i^2\}$
% \begin{enumerate}
	% \item \Hoare{s+(2*i-1)=i^2}{s:=s+(2*i-1)}{s=i^2} \hfill (Assignment)
	% \item
	% \item \Hoare{s=i^2 \wedge i \neq n}{i:=i+1}{s+(2*i-1)=i^2}
	% \item \Hoare{s = i^2 \wedge i \neq n}{i:=i+1;\ s:=s+(2*i-1)}{s=i^2} \hfill (1,3, Sequencing)
% \end{enumerate}
% \vspace{.65cm}
% \end{example}
% \end{frame}
% -------------------------------------------------------
% \addtocounter{framenumber}{-1}
% -------------------------------------------------------
% \begin{frame}{A simple program}
% \begin{example}[cont.]
% Check $P \equiv (s = i^2)$ is an invariant: \intens{prove \Hoare{P \wedge (i \neq n)}{S}{P}}
%
% \begin{center}
% $\infer
% {\HoareM{s = i^2 \wedge i \neq n}{i:=i+1;\ s:=s+(2*i-1)}{s=i^2}}
% {\HoareM{s=i^2 \wedge i \neq n}{i:=i+1}{Q} & \HoareM{Q}{s:=s+(2*i-1)}{s=i^2}}$
% \end{center}
%
% $Q$ is $\{s+(2*i-1)=i^2\}$
% \begin{enumerate}
	% \item \Hoare{s+(2*i-1)=i^2}{s:=s+(2*i-1)}{s=i^2} \hfill (Assignment)
        %  \item \Hoare{s+(2*(i+1)-1)=(i+1)^2}{i:=i+1}{s+(2*i-1)=i^2} \hfill  (Assignment)
	% \item \Hoare{s=i^2 \wedge i \neq n}{i:=i+1}{s+(2*i-1)=i^2} \onslide<2->{\hfill (2, Strength. Precond.)}
	% \item \Hoare{s = i^2 \wedge i \neq n}{i:=i+1;\ s:=s+(2*i-1)}{s=i^2} \hfill (1,3, Sequencing)
% \end{enumerate}
% \onslide<3->{So far, so good. }
% \end{example}
% \end{frame}
%-------------------------------------------------------

%-------------------------------------------------------
\begin{frame}{Un program simplu}
\begin{example}[cont.]
\intens{Completarea demonstrației \Hoare{\top}{Program}{s = n^2}}

\ldots
\begin{enumerate}[<+->]
{\setcounter{enumi}{3}}
	\item \Hoare{s = i^2 \wedge i \neq n}{i:=i+1;\ s:=s+(2*i-1)}{s=i^2}
	\item \Hoare{s=i^2}{while ... s:=s+(2*i-1)}{s=i^2 \wedge \neg (i \neq n)} \hfill(4, While)
	\item \Hoare{s=i^2}{while ... s:=s+(2*i-1)}{s=n^2}\\ \hfill(5, Slăbirea post-condiției)

	deoarece $s=i^2 \wedge \neg (i \neq n) \to s = n^2$
	\item \Hoare{0 = i^2}{s:= 0}{s = i^2} \hfill (Atribuire)
	\item \Hoare{0 = 0^2}{i:= 0}{0 = i^2} \hfill (Atribuire)
	\item \Hoare{\top}{i:= 0}{0 = i^2} \hfill (8, Întărirea pre-condiției)

	deoarece $\top \to 0 = 0^ 2$
	\item \Hoare{\top}{i := 0;\ s := 0}{s = i^2} \hfill(7,9, Secvențiere)
	\item \Hoare{\top}{Program}{s=n^2} \hfill(10, 6, Secvențiere)
\end{enumerate}
\end{example}
\end{frame}
%-------------------------------------------------------

%-------------------------------------------------------
%\begin{frame}{Hoare Logic}
%
%\begin{center}
%\begin{tabular}{ccc}
%\alert{Syntax} & \alert{Semantics} & \alert{Calculus} \\ \hline
%
%$\neg \wedge \vee \to \forall \exists$ & FOL &  N/A\\ \hline
%
%$= + - \geq \leq \mydots$ & Arithmetics &  N/A\\ \hline
%
%\intens{\texttt{:= ; while}}  & State maps variables & N/A \\
%\intens{\texttt{if then else}} & to values (no pointers) &  \\ \hline
%
%\alert{$\{P\}$\texttt{ S }$\{Q\}$} &  If initial state satisfies $P$ & 6 Inference\\
%& and \texttt{S} terminates then & Rules \\
%& final state satisfies $Q$&
%
%\end{tabular}
%\end{center}
%\end{frame}
%-------------------------------------------------------

%-------------------------------------------------------
\begin{frame}{Bibliografie}

{\footnotesize
\begin{itemize}
	\item Lecture Notes on "Formal Methods for Software Engineering", Australian National University, Rajeev Gor\'{e}.
	\item Mike Gordon, "Specification and Verification I", chapters 1 and 2.
	\item Michael Huth, Mark Ryan, "Logic in Computer Science: Modeling and Reasoning about Systems", 2nd edition, Cambridge University Press, 2004.
	\item Krzysztof R. Apt, Frank S. de Boer, Ernst-R\"{u}diger Olderog, "Verification of Sequential and Concurrent Programs", 3rd edition,  Springer.
\end{itemize}
}
\end{frame}
%-------------------------------------------------------
